% ==============================================================================
%  5.4  Inferencias sobre los parámetros
%  5.5  Inferencias sobre la respuesta media y predicción
% ==============================================================================
\section{Inferencias sobre los parámetros de la regresión}
\label{sec:t5-parametros}

La distribución de $\bhat=\XtXi\vect{X}\T\vect{Y}$ se obtiene a partir del siguiente resultado
sobre la normal multivariante.

\begin{teorema}[Transformaciones lineales de la normal multivariante][teo:t5-normal-lineal]
  Sea $\vect{U}\sim\Normal(\vect{\mu},\vect{\Sigma})$ un vector normal multivariante y
  $\vect{V}=\vect{c}+\vect{D}\vect{U}$ una transformación lineal de $\vect{U}$, donde $\vect{c}$ es
  un vector y $\vect{D}$ una matriz. Entonces
  \[
    \vect{V}\sim\Normal\big(\vect{c}+\vect{D}\vect{\mu},\ \vect{D}\vect{\Sigma}\vect{D}\T\big).
  \]
\end{teorema}

\begin{teorema}[Distribución de $\bhat$][teo:t5-dist-beta]
  \[
    \bhat\sim\Normal\big(\vect{\beta},\ \sigma^2\XtXi\big).
  \]
\end{teorema}

\begin{demostracion}
  Se tiene $\vect{Y}\sim\Normal(\vect{X}\vect{\beta},\sigma^2\vect{I})$ y, tomando
  $\vect{D}=\XtXi\vect{X}\T$, $\bhat=\vect{D}\vect{Y}$ es una transformación lineal de $\vect{Y}$.
  Por el \cref{teo:t5-normal-lineal}, $\bhat$ es normal, con
  \begin{align*}
    \E(\bhat) &= \vect{D}\vect{X}\vect{\beta}
      = \underbrace{\XtXi\vect{X}\T\vect{X}}_{\vect{I}}\vect{\beta} = \vect{\beta}, \\
    \Var(\bhat) &= \vect{D}\,\sigma^2\vect{I}\,\vect{D}\T
      = \sigma^2\big(\XtXi\vect{X}\T\big)\big(\XtXi\vect{X}\T\big)\T \\
      &= \sigma^2\XtXi\underbrace{\XtX\big(\XtXi\big)\T}_{\vect{I}} = \sigma^2\XtXi ,
  \end{align*}
  ya que $\XtX$ y su inversa son simétricas: $\big(\XtXi\big)\T=\XtXi$ y
  $\XtX\big(\XtXi\big)\T=\vect{I}$.
\end{demostracion}

Como $MSE$ es un estimador insesgado de $\sigma^2$ (\cref{prop:t2-mse}), la matriz de
varianzas-covarianzas de $\bhat$ se estima por
\[
  \widehat{\Var}(\bhat) = MSE\,\XtXi .
\]

\begin{proposicion}[Varianzas en la regresión simple][prop:t5-var-simple]
  Con la expresión de $\XtXi$ de la \cref{prop:t5-simple},
  \[
    \Var(\bhat) = \sigma^2\begin{pmatrix}
      \dfrac1n+\dfrac{\media{X}^2}{\SSX} & -\dfrac{\media{X}}{\SSX}\\[10pt]
      -\dfrac{\media{X}}{\SSX} & \dfrac{1}{\SSX}
    \end{pmatrix}:
  \]
  la diagonal contiene las varianzas de $\est{\beta}_0$ y $\est{\beta}_1$ del Tema 3
  (\cref{teo:t3-dist-b0,teo:t3-dist-b1}), y fuera de ella aparece su covarianza,
  $\Cov(\est{\beta}_0,\est{\beta}_1)=-\sigma^2\media{X}/\SSX$.
\end{proposicion}

\section{Inferencias sobre la respuesta media y predicción}
\label{sec:t5-respuesta-media}

\subsection{Respuesta media}

Dado el vector $\vect{X}_h=(1,\ x_h)\T$, la respuesta media en $\vect{X}_h$ es
\[
  \E(Y_h) = \E(Y\cond \vect{X}_h) = \vect{X}_h\T\vect{\beta},
\]
y su estimador es
\[
  \est{Y}_h = \vect{X}_h\T\bhat .
\]
$\est{Y}_h$ estima el parámetro $\E(Y_h)=\vect{X}_h\T\vect{\beta}=\beta_0+\beta_1x_h$, la media de la
distribución de $Y$ condicionada a $\vect{X}_h$, y no el valor de una observación concreta
(\cref{sec:t3-respuesta-media}). La estimación tiene sentido siempre que $\vect{X}_h$ tome valores
que podrían haberse dado en la muestra considerada, es decir, dentro del \emph{scope} del modelo:
la región de valores cubierta conjuntamente por todas las variables independientes.

\begin{proposicion}[Distribución de $\est{Y}_h$][prop:t5-Yh]
  Bajo los supuestos del modelo, el estimador $\est{Y}_h=\vect{X}_h\T\bhat$ de $\E(Y_h)$ es:
  \begin{itemize}
    \item normal, puesto que $\bhat$ lo es;
    \item insesgado: $\E(\est{Y}_h)=\E(\vect{X}_h\T\bhat)=\vect{X}_h\T\vect{\beta}=\E(Y_h)$;
    \item de varianza
      $\Var(\est{Y}_h)=\vect{X}_h\T\Var(\bhat)\vect{X}_h=\sigma^2\vect{X}_h\T\XtXi\vect{X}_h$.
  \end{itemize}
  La varianza estimada es
  \[
    \widehat{\Var}(\est{Y}_h) = \vect{X}_h\T\widehat{\Var}(\bhat)\vect{X}_h
    = MSE\cdot\vect{X}_h\T\XtXi\vect{X}_h .
  \]
\end{proposicion}

En la regresión simple, el mismo cálculo que el de $h_{ik}$ en la \cref{prop:t5-simple} da
\[
  \vect{X}_h\T\XtXi\vect{X}_h=\frac1n+\frac{(x_h-\media{X})^2}{\SSX},
\]
que es la expresión del \cref{teo:t3-dist-Yh}.

\subsection{Predicción de una nueva observación}

El objetivo es predecir una nueva observación de $Y$, $\Ynew$, cuando se conocen los valores de
la variable independiente, $\vect{X}_h$.
\begin{itemize}
  \item La nueva observación puede verse como el resultado de una nueva repetición del
    experimento, independiente de los resultados anteriores en los que se basa el análisis de
    regresión.
  \item Se asume que el modelo de regresión subyacente sigue siendo apropiado para la nueva
    observación.
\end{itemize}

Para hacer inferencias sobre $\Ynew$ es necesario conocer un estimador y la distribución del
error que se comete (\cref{sec:t3-prediccion}).
\begin{itemize}
  \item El estimador puntual sigue siendo $\est{Y}_h=\vect{X}_h\T\bhat$.
  \item Al estimar la respuesta media $\E(Y_h)$ se estima la media de la distribución de $Y$; al
    predecir una nueva respuesta $\Ynew$ se quiere predecir un resultado individual de la
    distribución de $Y$, que se situará alrededor de su media.
  \item Hay que tener en cuenta dos fuentes de variación: la variación en la posible ubicación
    del centro de la distribución de $Y$ y la variación dentro de la distribución de
    probabilidad de $Y$.
\end{itemize}
Como $\Ynew$ es independiente de la muestra, y por tanto de $\est{Y}_h$, la varianza del error de
predicción es
\begin{align*}
  \Var(\Ynew-\est{Y}_h) &= \Var(\Ynew)+\Var(\est{Y}_h)
  = \sigma^2+\sigma^2\vect{X}_h\T\XtXi\vect{X}_h \\
  &= \sigma^2\big(1+\vect{X}_h\T\XtXi\vect{X}_h\big),
\end{align*}
con estimador
\[
  \widehat{\Var}(\Ynew-\est{Y}_h) = MSE\big(1+\vect{X}_h\T\XtXi\vect{X}_h\big).
\]

\paragraph{Media de $m$ nuevas observaciones.} Para predecir la media $\Ymnew$ de $m$ nuevas
observaciones en $\vect{X}_h$: la varianza asociada a una observación es $\sigma^2$ y la asociada a
la media de $m$ observaciones es $\sigma^2/m$. La varianza del error de predicción es la suma de la
varianza de esa media y la de la estimación de la respuesta media:
\[
  \Var(\Ymnew-\est{Y}_h) = \sigma^2\left(\frac1m+\vect{X}_h\T\XtXi\vect{X}_h\right).
\]

\subsection{Intervalos de confianza y de predicción}

Sustituyendo $\sigma^2$ por $MSE$, los estadísticos tipificados siguen una $\tdist{n-2}$ y se
obtienen, en forma matricial, los mismos intervalos y contrastes del Tema 3.

\begin{resumen}[Inferencias en forma matricial (nivel $1-\alpha$)]
  Sea $c_{jj}$, $j=0,1$, el elemento $j$-ésimo de la diagonal de $\XtXi$.
  \begin{itemize}
    \item Coeficientes:
      $\displaystyle \beta_j\in\left(\est{\beta}_j\pm\tq{n-2}{1-\alpha/2}\sqrt{MSE\,c_{jj}}\right)$;
      para $H_0\colon\beta_j=b_j$,
      $\displaystyle t^*=\frac{\est{\beta}_j-b_j}{\sqrt{MSE\,c_{jj}}}\bajoHo\tdist{n-2}$.
    \item Respuesta media:
      $\displaystyle \E(Y_h)\in\left(\est{Y}_h\pm\tq{n-2}{1-\alpha/2}
      \sqrt{MSE\cdot\vect{X}_h\T\XtXi\vect{X}_h}\right)$.
    \item Nueva observación:
      $\displaystyle \Ynew\in\left(\est{Y}_h\pm\tq{n-2}{1-\alpha/2}
      \sqrt{MSE\big(1+\vect{X}_h\T\XtXi\vect{X}_h\big)}\right)$.
    \item Media de $m$ nuevas observaciones:
      $\displaystyle \Ymnew\in\left(\est{Y}_h\pm\tq{n-2}{1-\alpha/2}
      \sqrt{MSE\left(\tfrac1m+\vect{X}_h\T\XtXi\vect{X}_h\right)}\right)$.
  \end{itemize}
\end{resumen}
